\input{configTD}

% ── Poly à trous : commandes spécifiques ──────────────────────────
\newcommand{\atrou}[2][3cm]{%
  \ifthenelse{\boolean{showSolutions}}{%
    \par\vspace{0.5em}%
    \begin{mdframed}[backgroundcolor=ThemeLight, linewidth=0pt,
      skipabove=0pt, skipbelow=0pt,
      innertopmargin=8pt, innerbottommargin=8pt]%
    #2%
    \end{mdframed}%
    \par\vspace{0.5em}%
  }{%
    \par\vspace{0.5em}%
    \noindent\fbox{\begin{minipage}[c][#1]%
      {\dimexpr\textwidth-2\fboxsep-2\fboxrule\relax}%
      \color{white}\rule{0pt}{0pt}%
    \end{minipage}}%
    \par\vspace{0.5em}%
  }%
}

\newcommand{\val}[1]{%
  \ifthenelse{\boolean{showSolutions}}{#1}{\phantom{#1}}%
}

\newcommand{\R}{\mathbb{R}}

\title{\sffamily\bfseries Probabilités --- CM3\\[0.3em]
  \Large Lois continues classiques : uniformes, exponentielles, \dots}
\author{}
\date{21 avril 2026}

\begin{document}
\sffamily
\maketitle
\thispagestyle{empty}

% ====================================================================
\section*{1. Un même chantier, plusieurs variables continues}
% ====================================================================

Les lois continues servent à modéliser des grandeurs \textbf{mesurées}
(longueurs, résistances, températures) et des \textbf{durées}
((temps d'attente), temps entre pannes, délais).

\begin{center}
\begin{mdframed}[backgroundcolor=ThemeLight, linecolor=Theme,
  linewidth=1.5pt, innertopmargin=10pt, innerbottommargin=10pt, nobreak=true]
\centering\large\itshape
\og On planifie un coulage.\\
Quelle est la probabilité qu'il faille attendre \textbf{plus de 2 jours}
avant une fenêtre météo favorable ?\fg{}
\end{mdframed}
\end{center}

Ce CM introduit des \textbf{modèles standards} (lois)
et des \textbf{réflexes de calcul}.

% ====================================================================
\section*{2. Rappels minimaux : densité, répartition, intégrales}
% ====================================================================

Si \(X\) a une densité \(f\), alors
\[
  P(a\le X\le b)=\int_a^b f(x)\,dx,\qquad F(x)=P(X\le x)=\int_{-\infty}^x f(t)\,dt.
\]
Et
\[
  E[X]=\int_{-\infty}^{+\infty} x f(x)\,dx,\qquad \text{Var}(X)=E[X^2]-E[X]^2.
\]

% ====================================================================
\section*{3. Loi uniforme \(\mathcal{U}([a,b])\) : \og aucune raison de favoriser \fg{}}
% ====================================================================

\textbf{Quand ?}
Quand on sait seulement que \(X\) est \textbf{bornée} entre \(a\) et \(b\)
et qu'on suppose toutes les valeurs \og également plausibles \fg{}.

\begin{mdframed}[backgroundcolor=ThemeLight, linecolor=Theme,
  linewidth=1pt, innertopmargin=8pt, innerbottommargin=8pt, nobreak=true]
\textbf{Définition.}
\(X\sim \mathcal{U}([a,b])\) si
\[
  \boxed{f(x)=\frac{1}{b-a}\,\mathbf{1}_{[a,b]}(x).}
\]
Alors
\[
  \boxed{E[X]=\frac{a+b}{2}},\qquad
  \boxed{\text{Var}(X)=\frac{(b-a)^2}{12}}.
\]
\end{mdframed}

\textbf{Exercice (tolérance).}
Une erreur de réglage \(X\) (mm) est supposée uniforme sur \([-2,2]\).
Calculer \(P(|X|\le 1)\).

\atrou[3.2cm]{%
\[
  P(|X|\le 1)=P(-1\le X\le 1)=\frac{1-(-1)}{2-(-2)}=\frac{2}{4}=\frac12.
\]
}

% ====================================================================
\section*{4. Loi exponentielle \(\mathcal{E}(\lambda)\) : \og attendre un événement \fg{}}
% ====================================================================

\textbf{Quand ?}
Pour modéliser un \textbf{temps d'attente} avant un événement
qui arrive \og au hasard \fg{} avec un taux moyen constant :
panne, incident, arrivée d'une opportunité, \dots

\begin{mdframed}[backgroundcolor=ThemeLight, linecolor=Theme,
  linewidth=1pt, innertopmargin=8pt, innerbottommargin=8pt, nobreak=true]
\textbf{Définition.}
Pour \(\lambda>0\), \(T\sim\mathcal{E}(\lambda)\) si
\[
  \boxed{f(t)=\lambda e^{-\lambda t}\,\mathbf{1}_{t\ge 0}.}
\]
Sa répartition vaut
\[
  \boxed{F(t)=P(T\le t)=1-e^{-\lambda t}\quad (t\ge 0).}
\]
Et
\[
  \boxed{E[T]=\frac{1}{\lambda}},\qquad
  \boxed{\text{Var}(T)=\frac{1}{\lambda^2}}.
\]
\end{mdframed}

\textbf{Probabilité de dépasser un délai.}
Pour \(t\ge 0\),
\[
  P(T>t)=1-F(t)=e^{-\lambda t}.
\]

\textbf{Exercice (fenêtre météo).}
On modélise le temps d'attente \(T\) (jours) avant un jour favorable par
\(\mathcal{E}(\lambda)\) avec \(E[T]=1\) jour.
Calculer \(P(T>2)\).

\atrou[3.8cm]{%
Ici \(E[T]=1/\lambda=1\), donc \(\lambda=1\).
\[
  P(T>2)=e^{-1\cdot 2}=e^{-2}\approx 0{,}1353.
\]
}

\textbf{Propriété sans mémoire (analogue continu de la géométrique).}
Pour \(s,t\ge 0\),
\[
  P(T>s+t\mid T>s)=P(T>t).
\]

\atrou[4.2cm]{%
\textbf{Preuve.}
\[
  P(T>s+t\mid T>s)=\frac{P(T>s+t)}{P(T>s)}
  =\frac{e^{-\lambda(s+t)}}{e^{-\lambda s}}=e^{-\lambda t}=P(T>t).
\]
}

% ====================================================================
\section*{5. Passage par la répartition : un réflexe de calcul}
% ====================================================================

Pour beaucoup de lois, \(F\) est plus simple que \(f\).
Deux exemples utiles :

\begin{enumerate}[leftmargin=1.2em]
  \item \(P(a\le X\le b)=F(b)-F(a)\).
  \item Si \(X\) est continue, \(P(X<a)=P(X\le a)\) (pas de masse en un point).
\end{enumerate}

% ====================================================================
\section*{6. Changement de variable simple : linéaire et monotone}
% ====================================================================

\textbf{Transformation affine.}
Si \(Y=aX+b\), alors
\[
  E[Y]=aE[X]+b,\qquad \text{Var}(Y)=a^2\text{Var}(X).
\]

\textbf{Exercice.}
Si \(T\) est en jours et \(H=24T\) en heures, exprimer \(E[H]\) et \(\text{Var}(H)\).

\atrou[3.4cm]{%
\[
  E[H]=24E[T],\qquad \text{Var}(H)=24^2\,\text{Var}(T).
\]
}

\textbf{Cas monotone (par la répartition).}
Si \(g\) est croissante et \(Y=g(X)\), alors
\[
  F_Y(y)=P(Y\le y)=P(g(X)\le y)=P(X\le g^{-1}(y))=F_X(g^{-1}(y)).
\]

\textbf{Exercice (max).}
Supposons \(T\sim\mathcal{E}(\lambda)\) et \(Y=\min(T,3)\) (on attend au maximum 3 jours).
Donner \(P(Y<3)\).

\atrou[3.6cm]{%
\[
  P(Y<3)=P(\min(T,3)<3)=P(T<3)=F_T(3)=1-e^{-3\lambda}.
\]
}

% ====================================================================
\section*{7. Où se place la loi normale ?}
% ====================================================================

La normale \(\mathcal{N}(\mu,\sigma^2)\) est typiquement un modèle de
\textbf{mesure} et d'\textbf{erreur}.
On la manipule surtout via la standardisation
\(\;Z=(X-\mu)/\sigma\sim\mathcal{N}(0,1)\).

\medskip

Dans la suite (CM4), on l'utilisera pour construire des
\textbf{intervalles de confiance} sur une moyenne.

% ====================================================================
\section*{8. Bilan : choisir un modèle}
% ====================================================================

\atrou[8cm]{%
\begin{enumerate}
  \item \textbf{Uniforme} \(\mathcal{U}([a,b])\) : variable bornée, pas de valeur privilégiée.
  \item \textbf{Exponentielle} \(\mathcal{E}(\lambda)\) : temps d'attente, taux constant,
    propriété sans mémoire.
  \item \textbf{Normale} \(\mathcal{N}(\mu,\sigma^2)\) : mesures/erreurs, cloche,
    calcul via standardisation.
  \item Pour calculer : passer par \(F\), utiliser \(P(a\le X\le b)=F(b)-F(a)\),
    et exploiter les transformations simples.
\end{enumerate}
}

\end{document}

